\chapter{Escenarios de despliegue}

La implantación se divide en tres escenarios para evitar un salto brusco. Cada escenario amplía el alcance operativo y el número de robots solo cuando los indicadores del paso anterior son aceptables.

\begin{longtable}{P{0.18\textwidth}P{0.25\textwidth}P{0.22\textwidth}P{0.25\textwidth}}
\toprule
Escenario & Objetivo y alcance & Riesgos & Beneficios y criterios de éxito\\
\midrule
\endhead
Piloto con 1 robot & Validar navegación, kitting básico, RFID/QR e inspección FOD en una zona reducida. Coste aproximado: 237.444 \texteuro{} con contingencia técnica y margen incluidos. & Rechazo de usuarios, mapa incompleto, interferencias en pasillos. & 95\% de misiones completadas, cero incidentes de seguridad, reducción medible de esperas en estaciones piloto.\\
Expansión con 7 robots & Cubrir varias estaciones HTP, retornos de útiles y rondas FOD programadas. Coste aproximado: 1.100.252 \texteuro{} con contingencia técnica y margen incluidos. & Congestión entre robots, prioridades mal configuradas, dependencia de red WiFi. & Entregas trazadas, menor búsqueda manual, gestor de flota estable y aceptación operativa.\\
Despliegue extendido con 17 robots & Soportar flujo ampliado hacia 80 HTP/mes o más, con coordinación multi-zona e integración avanzada. Coste aproximado: 2.495.768 \texteuro{} con contingencia técnica y margen incluidos. & Complejidad de mantenimiento, necesidad de gobierno de datos, integración MES/ERP más exigente. & Operación multi-robot estable, registros completos, tiempos improductivos reducidos y capacidad de escalado medida.\\
\bottomrule
\end{longtable}

\section{Criterios de paso entre escenarios}

El paso de 1 a 7 robots requiere que el piloto demuestre seguridad, fiabilidad de misión y utilidad para los operarios. El paso de 7 a 17 robots requiere un gestor de flota probado, mantenimiento preventivo definido, red industrial estable y reglas de prioridad acordadas con producción, calidad y seguridad.
